% =========================================================
% Perturbative Operator Framework
% =========================================================

\section{Perturbative Operator Framework}

\subsection{Perturbative Structure}

We consider perturbations of the representative background geometry introduced in the previous section.

The metric perturbation is written as

\begin{equation}
g_{ab}
\rightarrow
g_{ab} + h_{ab},
\end{equation}

where the perturbative sector satisfies

\begin{equation}
|h_{ab}| \ll 1.
\end{equation}

The present analysis is restricted to the linear perturbative regime.

No nonlinear stability claims are made.

% =========================================================

\subsection{Spherical Symmetry and Mode Reduction}

Due to the static spherical symmetry of the representative background geometry, perturbative modes may be decomposed into angular sectors.

The perturbative analysis is formulated using a reduced radial description associated with angular harmonic decomposition.

The present work does not attempt a complete tensor-harmonic classification. Instead, we focus on the resulting effective radial perturbative structure.

This reduced framework is interpreted as a representative operator-theoretic sector suitable for perturbative admissibility analysis.

% =========================================================

\subsection{Tortoise Coordinate}

We introduce the tortoise coordinate $r_*$ through

\begin{equation}
\frac{dr_*}{dr}
=
\frac{1}{F(r)},
\end{equation}

where $F(r)$ is the representative lapse function defined in Eq.~(\ref{eq:lapse}).

The tortoise coordinate regularizes the perturbative radial structure near simple horizon sectors and allows the perturbation equations to be expressed in Schrödinger-type form.

% =========================================================

\subsection{Effective Radial Operator}

The perturbative radial dynamics are modeled formally through the operator

\begin{equation}
L_{\mathrm{TIG}}
=
-
\frac{d^2}{dr_*^2}
+
V_{\mathrm{eff}}(r).
\label{eq:operator}
\end{equation}

The effective potential is represented by the candidate structure

\begin{equation}
V_{\mathrm{eff}}(r)
=
F(r)
\left[
\frac{\ell(\ell+1)}{r^2}
+
\frac{F'(r)}{r}
\right].
\label{eq:potential}
\end{equation}

Here $\ell$ denotes the angular harmonic index.

The exact perturbative derivation of the effective potential remains incomplete and is treated conservatively throughout the present work.

Accordingly, Eq.~(\ref{eq:potential}) should be interpreted as a representative candidate perturbative potential compatible with the analyzed geometric sector.

% =========================================================

\subsection{Asymptotic Behavior of the Potential}

For large radial distance,

\begin{equation}
r \gg r_c,
\end{equation}

the effective potential approaches the Schwarzschild-type asymptotic structure

\begin{equation}
V_{\mathrm{eff}}(r)
\sim
\frac{\ell(\ell+1)}{r^2}
+
\mathcal{O}\!\left(\frac{1}{r^3}\right).
\end{equation}

Thus the perturbative sector remains compatible with the classical far-field regime.

Near the regulated core, the bounded behavior of $F(r)$ prevents the emergence of divergent representative potential terms within the analyzed perturbative sector.

% =========================================================

\subsection{Quadratic Form Structure}

Associated with the formal operator (\ref{eq:operator}), we define the quadratic form

\begin{equation}
Q[\psi]
=
\int
\left(
\left|
\frac{d\psi}{dr_*}
\right|^2
+
V_{\mathrm{eff}}(r)
|\psi|^2
\right)
dr_*.
\label{eq:quadratic_form}
\end{equation}

The quadratic form plays a central role in the perturbative admissibility structure.

In particular, lower-boundedness of $Q[\psi]$ is required for controlled spectral behavior and bounded linear evolution.

The present work does not establish a complete positivity theorem for the quadratic form.

Such results remain conditional upon further operator-theoretic analysis.

% =========================================================

\subsection{Self-Adjoint Realization}

The perturbative framework requires a mathematically controlled realization of the formal operator $L_{\mathrm{TIG}}$.

The present analysis adopts a conservative operator-theoretic perspective in which self-adjoint realization is treated conditionally.

Formally, one seeks a self-adjoint extension

\begin{equation}
L_{\mathrm{TIG}}^{F},
\end{equation}

obtained through quadratic-form methods and admissible endpoint classification.

The existence and uniqueness of such a realization depend on:
\begin{itemize}[leftmargin=0.6cm]
\item endpoint behavior,
\item domain structure,
\item quadratic-form closability,
\item and lower-boundedness conditions.
\end{itemize}

The complete proof of self-adjoint realization remains open.

% =========================================================

\subsection{Spectral Admissibility}

The central spectral admissibility condition of the present framework is

\begin{equation}
\sigma(L_{\mathrm{TIG}}^{F})
\subset
[0,\infty),
\label{eq:spectral_condition}
\end{equation}

where $\sigma(L_{\mathrm{TIG}}^{F})$ denotes the spectrum of the self-adjoint perturbative operator.

Condition (\ref{eq:spectral_condition}) excludes negative spectral modes corresponding to exponentially growing perturbative instabilities.

The present work does not establish a rigorous spectral nonnegativity theorem.

Instead, Eq.~(\ref{eq:spectral_condition}) defines the conditional spectral admissibility requirement underlying the TIG3 perturbative program.

% =========================================================

\subsection{Conditional Bounded Linear Evolution}

Under the assumptions:
\begin{itemize}[leftmargin=0.6cm]
\item admissible endpoint structure,
\item self-adjoint realization,
\item and spectral nonnegativity,
\end{itemize}

the perturbative evolution equation

\begin{equation}
\partial_t^2 \Psi
+
L_{\mathrm{TIG}}^{F}\Psi
=
0
\end{equation}

admits bounded linear finite-energy evolution within the analyzed perturbative sector.

This statement is conditional.

No nonlinear boundedness theorem is established in the present work.

% =========================================================

\subsection{Near-Critical Perturbative Sector}

Near the critical branch-merger regime

\begin{equation}
\beta \approx \beta_c,
\end{equation}

the perturbative geometry may develop modified spectral structures associated with elongated throat behavior.

Potential effects include:
\begin{itemize}[leftmargin=0.6cm]
\item slow decay,
\item metastable trapping,
\item resonance accumulation,
\item and spectral compression.
\end{itemize}

The present work does not establish rigorous resonance theorems in the near-critical sector.

Such structures remain conjectural and require future asymptotic spectral analysis.

% =========================================================

\subsection{Interpretation}

The perturbative framework developed above should be interpreted as a conditional operator-theoretic admissibility program associated with a representative bounded-curvature vacuum geometry.

The present analysis establishes:
\begin{itemize}[leftmargin=0.6cm]
\item a coherent perturbative radial operator structure,
\item a quadratic-form framework,
\item and a conditional spectral admissibility hierarchy.
\end{itemize}

The framework does not currently establish:
\begin{itemize}[leftmargin=0.6cm]
\item nonlinear stability,
\item complete spectral classification,
\item exact perturbative derivation of the effective potential,
\item or full covariant gravitational closure.
\end{itemize}

The present work is therefore interpreted conservatively as a mathematically structured perturbative admissibility framework suitable for further operator-theoretic development.
